\chapter*{解题复盘卡}\addcontentsline{toc}{chapter}{解题复盘卡}
\begin{insight}{用一句话说出本题的结构}
“两列只差一个简单向量”“非齐次解集是核的平移”“重根还要算几何重数”“配方是在做合同变换”。能够给出这种具体判断，再选工具，通常比一上来展开全部元素更有效。
\end{insight}
\begin{enumerate}
\item \textbf{先核对对象。}矩阵是否方形，向量属于哪个空间，未知量与方程各有多少个？
\item \textbf{列出有效条件。}可逆、对称、正定、无关、满秩分别在哪一步用到？不要把题中未给的条件补进证明。
\item \textbf{画一个可解释的小图。}用二维或三维模型看独立方向、约束交集、投影和主轴；注明这是模型，再用代数推广。
\item \textbf{寻找结构捷径。}低秩找外积，同背景找更新，多项式关系先降次，特殊子空间先分解，参数问题先查退化值。
\item \textbf{完成必要输出。}求通解要写全部自由参数，求变换要给出可逆或正交矩阵，求极值要写等号条件，分类要覆盖每个参数分支。
\item \textbf{用另一种语言核验。}计算后看秩与维数、迹与行列式、符号与惯性、残差与正交是否一致。
\end{enumerate}
\begin{warning}
最短的计算不一定最适合第一次学习。第一次应理解定义与条件，第二次比较方法，第三次再练习限时选择。发现结果与图形直觉冲突时，检查条件与计算，而不是凭图形否定一个已经完整证明的结论。
\end{warning}
